% Role C contribution included by main.tex.
% Source: Zhou et al., Sections 4, 4.1--4.3, Figure 5, pp. 324--327.
% Uses bibliography key zhou2003consistency from references.bib.
\subsection{Experimental evidence and its limits}
Zhou et al.~\cite{zhou2003consistency} evaluate their method using a
two-moons illustration, a four-class subset of USPS handwritten digits
(3,874 examples), and four newsgroup topics (3,970 documents represented
by 8,014 features). The real-data comparisons include 1-nearest neighbor,
one-versus-rest RBF SVMs, a harmonic Gaussian field method with class mass
normalization, and two consistency variants. Error is measured on unrevealed
labels, with the digit and text results averaged over 100 trials and at least
one labeled example per class. The consistency methods use $\alpha=0.99$.

Figure~5 of the original paper shows substantially lower digit-recognition
error for the main consistency method than for the supervised baselines.
The text results are more qualified: the harmonic method performs better
at some label budgets. These findings support exploiting unlabeled structure
under the tested conditions, rather than universal superiority over other
semi-supervised methods. The authors' explanation of the harmonic method's
nonmonotonic text results emphasizes the reliability of class-proportion
estimates used in class mass normalization.

A central limitation is model selection. The authors use each algorithm's
respective optimal parameters while acknowledging the difficulty of choosing
parameters with very few labels. The reported comparisons therefore leave
practical label-scarce tuning unresolved. In addition, ensuring label coverage
of every class is a favorable assumption, and the absence of displayed error
bars in Figure~5 limits interpretation of small differences between curves.
These are our assessments of the evidence, not additional experimental findings.

\subsection{Our partial reproduction on two moons}
For each of ten fixed seeds (0--9), we generated 400 two-moons points
with Gaussian coordinate noise of standard deviation 0.08. We revealed one
uniformly sampled label per class and scored predictions on the other 398
points. Euclidean 1-NN and our NumPy LGC implementation used the same
coordinates and revealed labels in each trial.

At our fixed setting of $\sigma=0.15$ and $\alpha=0.99$, mean error across
seeds was $0.45\%$ for LGC (sample SD 0.54 percentage points), compared with
$22.36\%$ for 1-NN (SD 4.22 percentage points). LGC had lower error in every
paired trial. Its per-seed accuracies also matched the independent fixed-point
check described in Section~\ref{sec:theory}.

A descriptive sweep held $\alpha$, data, and revealed labels fixed. LGC mean
error was $0.15\%$, $0.10\%$, $0.45\%$, $14.07\%$, and $32.59\%$ at
$\sigma=0.05, 0.10, 0.15, 0.30,$ and $0.60$, respectively. The results show
that label propagation can exploit the curved structure, and that the
advantage depends on the graph bandwidth. The sweep describes sensitivity
on these trials. Our results concern generated toy data with a label from
each class, so their scope is narrower than the paper's image and text studies.
